\clearpage\pagestyle{matSimple}\phantomsection\label{sources}\pdfbookmark[0]{Sources}{bm-sources}
\matPageTitre{Sources et édition}{Références officielles · France}
{\fontsize{9.5}{13.5}\selectfont \textbf{1. Programme d’enseignement de l’école maternelle}\par Arrêté publié au BO n° 41 du 31 octobre 2024, applicable à la rentrée 2025 ; version consolidée intégrant le BO n° 19 du 7 mai 2026. Domaine « L’acquisition des premiers outils mathématiques » et partie « Se repérer dans l’espace ».\par\href{https://eduscol.education.gouv.fr/sites/default/files/document/programme-cycle-1-consolide-127565.pdf}{\textcolor{otGreenDark}{Consulter la source officielle}}\par\medskip
\textbf{2. Ressources d’accompagnement Éduscol}\par Ressources pour le cycle 1 : construction du nombre, résolution de problèmes, formes et grandeurs, motifs organisés.\par\href{https://eduscol.education.gouv.fr/}{\textcolor{otGreenDark}{Consulter la source officielle}}\par\medskip
\textbf{3. Choix faits dans ce cahier}\par Les quantités à dénombrer vont jusqu’à dix ; une activité de parcours en compte quatorze, et la comptine va jusqu’à trente. Les nombres écrits ne dépassent pas dix, sauf dans les activités de comptine, où l’adulte aide à lire. Les problèmes comportent des énoncés « discordants » (« de moins » quand il faut ajouter), comme le demande le programme. Les solides ne sont jamais dessinés en perspective : on les décrit par leurs faces, dessinées à plat. Les assemblages comportent au plus huit pièces. Le nombre zéro n’est pas étudié. Chaque fiche commence par une activité avec de vrais objets, décrite dans le guide de l’adulte.\par}
